\clearpage\pagestyle{otMain}\setlength{\otBlocEspacement}{2.5mm}\fontsize{11}{15}\selectfont\corrigesChapitre{12}{Horaires et durées}
\begin{blocLivre}{cor-12-1}\begin{corrige}{1}{Lire l’heure}\demarche{Lire la petite aiguille, puis la grande.}
$\res{15\text{ h }20}$ ; $\res{19\text{ h }45}$ ; $\res{23\text{ h }05}$ : on ajoute 12 h aux heures lues sur le cadran.
\label{fig-cor-12-1}\figureGeo[corrige]{dur-1}
\end{corrige}\end{blocLivre}
\begin{blocLivre}{cor-12-2}\begin{corrige}{2}{Unités de durée}\demarche{Utiliser les relations entre unités.}
$\res{3\,600}$ s ; $\res{72}$ h ; $\res{135}$ s ; $\res{10}$ siècles.
\end{corrige}\end{blocLivre}
\begin{blocLivre}{cor-12-3}\begin{corrige}{3}{Fractions d’heure}\demarche{Prendre la fraction de 60 minutes.}
$\res{15}$ min ; $\res{45}$ min ; $\res{30}$ min ; $60:3=\res{20}$ min.
\end{corrige}\end{blocLivre}
\begin{blocLivre}{cor-12-4}\begin{corrige}{4}{Calculer une durée}\demarche{Avancer jusqu’aux heures rondes.}
9 h 47 → 10 h : 13 min ; 10 h → 11 h : 1 h ; 11 h → 11 h 15 : 15 min. Durée : $\res{1\text{ h }28\text{ min}}$.
\end{corrige}\end{blocLivre}
\begin{blocLivre}{cor-12-5}\begin{corrige}{5}{Heure d’arrivée}\demarche{Ajouter les heures, puis les minutes en passant par l’heure ronde.}
14 h 35 + 2 h = 16 h 35 ; + 25 min = 17 h ; + 25 min = 17 h 25. Arrivée : $\res{17\text{ h }25}$.
\end{corrige}\end{blocLivre}
\begin{blocLivre}{cor-12-6}\begin{corrige}{6}{Convertir}\demarche{Division euclidienne par 60.}
$200=60\times3+20$ : $\res{3\text{ h }20\text{ min}}$. $135=60\times2+15$ : $\res{2\text{ h }15\text{ min}}$. $1\,000=60\times16+40$ : $\res{16\text{ min }40\text{ s}}$.
\end{corrige}\end{blocLivre}
\begin{blocLivre}{cor-12-7}\begin{corrige}{7}{Heures décimales}\demarche{Multiplier la partie décimale par 60.}
$\res{1\text{ h }30\text{ min}}$ ; $\res{2\text{ h }15\text{ min}}$ ; $\res{6\text{ min}}$ ; $\res{3\text{ h }45\text{ min}}$.
\end{corrige}\end{blocLivre}
\begin{blocLivre}{cor-12-8}\begin{corrige}{8}{Le piège de la virgule}\demarche{Convertir la partie décimale en minutes.}
$\res{\text{Non}}$ : $0{,}3\times60=18$, donc $1{,}3$ h $=1$ h 18 min. Et 1 h 30 min $=\res{1{,}5}$ h, car 30 min est la moitié d’une heure.
\end{corrige}\end{blocLivre}
\begin{blocLivre}{cor-12-9}\begin{corrige}{9}{La matinée au collège}\demarche{Lire les instants sur la frise, puis calculer les durées.}
De 8 h 25 à 11 h 55 : $\res{3\text{ h }30\text{ min}}$. Récréation : de 10 h 05 à 10 h 20, $\res{15\text{ min}}$. Cours : $\res{3\text{ h }15\text{ min}}$.
\label{fig-cor-12-9}\figureGeo[corrige]{dur-9}
\end{corrige}\end{blocLivre}
\begin{blocLivre}{cor-12-10}\begin{corrige}{10}{Années bissextiles}\demarche{Tester la division par 4.}
$2\,028=4\times507$ : 2028 sera $\res{\text{bissextile}}$. 2027 compte 365 jours et 2028 en compte 366 : $\res{731}$ jours.
\end{corrige}\end{blocLivre}
\begin{blocLivre}{cor-12-11}\begin{corrige}{11}{Le train}\demarche{Calculer chaque durée, puis convertir 100 min.}
Paris-Lyon : $\res{1\text{ h }56\text{ min}}$. Lyon-Marseille : $\res{1\text{ h }44\text{ min}}$. Total : 1 h 56 min + 1 h 44 min $=2$ h 100 min $=\res{3\text{ h }40\text{ min}}$.
\end{corrige}\end{blocLivre}
\begin{blocLivre}{cor-12-12}\begin{corrige}{12}{La course}\demarche{Additionner chaque unité, puis convertir les 60.}
40 s + 35 s $=75$ s $=1$ min 15 s. 25 min + 12 min + 1 min $=38$ min. Temps : $\res{3\text{ h }38\text{ min }15\text{ s}}$.
\end{corrige}\end{blocLivre}
\begin{blocLivre}{cor-12-13}\begin{corrige}{13}{Le bus de ville}\demarche{Tout convertir en secondes.}
2 min 45 s $=165$ s. $12\times165=1\,980$ s $=\res{33\text{ min}}$, car $1\,980=60\times33$.
\end{corrige}\end{blocLivre}
\begin{blocLivre}{cor-12-14}\begin{corrige}{14}{Les calendriers}\demarche{Convertir un jour en minutes et comparer.}
Un jour $=1\,440$ min ; $1\,440=11\times130+10$ : environ $\res{131}$ ans par jour. De 325 à 1582 : $1\,257$ ans, soit $1\,257\times11=13\,827$ min, presque $\res{10}$ jours ($10\times1\,440=14\,400$). On a supprimé ces jours pour recaler le calendrier sur les saisons.
\end{corrige}\end{blocLivre}
\begin{blocLivre}{cor-12-15}\begin{corrige}{15}{Le salaire de la journée}\demarche{Passer en heures décimales pour multiplier.}
7 h 45 min $=7{,}75$ h, car 45 min font $\frac34$ d’heure. Salaire : $7{,}75\times14=\res{108{,}50}$ €.
\end{corrige}\end{blocLivre}
